\chapter{Технологическая часть}
В данной части будут рассмотрены средства реализации программного обеспечения, предоставлены листинги и приведены тестовын данные

\section{Средства реализации}
В качестве средства реализации был выбран язык C++~\cite{lit5}, из-за возможности управления памятью и наличия контейнерных классов, встроенных в стандартную библиотеку . Разработка проводилась в среде программирования VS Code. Для сборки использована утилита CMake. Для визуализации данных, полученных в ходе исследования использовался язык Python~\cite{lit6} за счёт обилия библиотек для анализа данных. Для построения графиков были использованы pandas~\cite{lit7} и matplotlib~\cite{lit8}. Для запуска тестов и последующего сохранения результатов также использовался язык Python.

\section{Реализация алгоритмов}
При выполнении лабораторной работы были реализованы два варианта алгоритма поиска максимума. Реализация итерационного алгоритма указан на листинге~\ref{list:iter}. Реализация рекурсивного алгоритма указан на листинге~\ref{list:rec}
\begin{lstlisting}[language=C++, caption={Итерационный алгоритм поиска максимума}, captionpos=b, label=list:iter]
	void Sequence::default_max(int &max) { \\ 0
		max = INT_MIN; \\ 1
		size_type i = 0; \\ 2
		while (i < _len) { \\ 3
			if (_container[i] > max) \\ 4
				max = _container[i];}} \\ 5
			i++; \\ 6
\end{lstlisting}

\begin{lstlisting}[language=C++, caption={Рекурсивный алгоритм поиска максимума}, captionpos=b,  label=list:rec]
	void Sequence::recursion_max(int &max, size_type index) { \\ 0
		if (index == _len) \\ 1
		return; \\ 2
		if (_container[index] > max) \\ 3
			max = _container[index]; \\ 4
		recursion_max(max, index += 1);} \\ 5
\end{lstlisting}



\section{Тестирование ПО}
Для тестирование реализованы ПО были выделены следующие тестовые случаи для каждого алгоритма:
\begin{table}[h]
	\centering
	\caption{Тестовые случаи для функции рекурсивного поиска максимума}
	\begin{tabular}{|c|c|c|c|}
		\hline
		\textbf{Входная последовательность} & \textbf{Ожидаемый максимум} &  \textbf{Описание теста} \\
		\hline
		\(\{-1,\ -3,\ 0\}\) & \(0\) & \text{Максимум ноль} \\
		\hline
		\(\{-1,\ -3,\ -6\}\) & \(-1\)  & \text{Максимум отрицательный}  \\
		\hline
		\(\{100,\ 23,\ 6\}\) & \(100\) & \text{Максимум положительный}  \\
		\hline
		\(\{100\}\)  & \(100\) & \text{Последовательность из 1 элемента} \\
		\hline
		\(\{100,\ 6\}\)  & \(100\) & \text{Последовательность из 2 элементов} \\
		\hline
		\(\{100,\ 50,\ 20,\ 40\}\) & \(100\) & \text{Максимум первый} \\
		\hline
		\(\{100,\ 50,\ 33,\ 200\}\)& \(200\) & \text{Максимум последний} \\
		\hline
		\(\{40,\ 20,\ 10,\ 0,\ 100,\ 50\}\) & \(100\) & \text{Максимум в середине}  \\
		\hline
	\end{tabular}
\end{table}

\section*{Вывод}
На основе конструкторской части были реализованы и протестированы алгоритмы нахождения максимума последовательности.

\clearpage
